% 第12章第4节：两种构建的等价关系（中间函数与接口消息的对应、树宽与最小诱导宽度、四环上的双向构建）。

\section{两种构建的等价关系}
\label{sec:graph-dp-construction-equivalence}

前两节分别从消元次序和树分解出发，构建出同一类结构。本节把两种表示逐项对照：消元簇对应袋，剩余作用域对应接口，中间函数对应向父袋传递的消息。比较规模时，要区分某一个构建的宽度与全部构建中的最优值；后者在两种表示下相等。最后在四环上把双向转换逐步做一遍。

\subsection{中间函数与接口消息的对应}
\label{sec:graph-dp-function-message-correspondence}

在算法~\ref{alg:graph-dp-order-to-decomposition}得到的袋树中，取消元森林的根为根。每个非根袋$B_i=\{v_i\}\cup S_i$与父袋的接口恰为$S_i$：第~\ref{sec:graph-dp-order-to-decomposition}小节已说明父袋包含$S_i$；另一方面，父袋若是$B_j=\{v_j\}\cup S_j$，其中的变量在处理$v_j$时都还在当前图中，而$v_i$此前已被删去，所以$v_i\notin B_j$，父袋若是根袋$B_R$，也因$v_i\notin R$而不含$v_i$。所以$B_i$的内部变量只有$v_i$一个，这与引理~\ref{lem:graph-dp-unique-internal}一致：$v_i$所在的最高节点就是$B_i$。

于是两种表示中的对象逐项对应：
\begin{itemize}
  \item 第$i$步的消元簇$B_i$对应袋树的一个节点；
  \item 中间函数$\psi_i$的作用域$S_i$对应该节点与父节点的接口；
  \item 消去$v_i$对应汇总该节点的内部变量；
  \item “$\psi_i$在消去$S_i$中最早的变量时被使用”对应“消息从子节点发往父节点”。
\end{itemize}
第~\ref{sec:graph-dp-elimination-message-correspondence}小节在相同的局部项归属下证明，按消元次序算出的$\psi_i$与沿袋树算出的消息逐项相等。

反过来，算法~\ref{alg:graph-dp-decomposition-to-order}从树分解读出消元次序时，处理叶袋$\ell$就是一次消息传递：消去$B_\ell\setminus B_p$中的变量，留下以$B_\ell\cap B_p$为作用域的结果交给父袋。一个袋可以有多个内部变量，相当于把连续的几步消元合并在一个计算单元里完成。

对应关系不要求逐表相同。把相邻的袋合并、把一个局部项从一个容纳它的袋移到另一个，都会改变中间表；只要结果仍是合法的树分解、每个局部项的作用域被所归的袋容纳、每项恰好归入一处，最终的最小成本、条件最小成本表、计数和加权总和都不变。这一点在第~\ref{chap:general-graph-dp-computation}章随各任务的正确性一起证明。

\paragraph{两种表示的组织差别}
两种做法的差别在于组织方式。变量消元边走边决定新函数交给谁，结构与数值交织在一起；树分解先把结构确定下来并核验，同一个分解可以用于作用域相同而数值不同的多组局部函数，也便于在不同的根下重复使用，例如第~\ref{sec:graph-dp-tree-two-pass}小节那样再发一遍反向消息，得到每个袋上的表。

\subsection{树宽与最小诱导宽度}
\label{sec:graph-dp-treewidth-induced-width}

定理~\ref{thm:graph-dp-elimination-decomposition}给出一个方向：每个消元次序都产生宽度等于其诱导宽度的树分解。第~\ref{sec:graph-dp-decomposition-to-order}小节给出另一个方向：每个宽度为$w$的树分解都读出诱导宽度不超过$w$的消元次序。两个方向合起来，得到下面的等式。

\begin{theorem}[树宽等于最小诱导宽度]
\label{thm:graph-dp-treewidth-induced-width}
设$\mathcal G$是非空有限无向图。则
\[
\operatorname{tw}(\mathcal G)=\min_{\pi}\ \max_{v}\ |N^+_\pi(v)|,
\]
其中$\pi$取遍全部顶点的消元次序，$N^+_\pi(v)$是按$\pi$消元时$v$的剩余邻居集合。
\end{theorem}

\begin{proof}
记右端为$w^\ast$，它是有限个次序上的最小值，因而可以达到。由定理~\ref{thm:graph-dp-elimination-decomposition}，诱导宽度为$w^\ast$的次序给出宽度为$w^\ast$的树分解，所以$\operatorname{tw}(\mathcal G)\leq w^\ast$。反过来，取宽度为$\operatorname{tw}(\mathcal G)$的树分解，由算法~\ref{alg:graph-dp-decomposition-to-order}及其正确性论证，得到诱导宽度不超过$\operatorname{tw}(\mathcal G)$的次序，所以$w^\ast\leq\operatorname{tw}(\mathcal G)$。
\end{proof}

这个等式说的是两种表示的最优值相等，不是两种对象一一对应。同一个树分解按不同的根、不同的叶袋处理次序可以读出许多消元次序；不同的消元次序也可能给出合并后相同的袋树。等式的用途是：寻找好的结构时，可以在消元次序上搜索，也可以在树分解上搜索，二者的最优规模相同。对某一个具体次序或分解，它的宽度只是树宽的上界。

\subsection{四环上的双向构建}
\label{sec:graph-dp-four-cycle-construction}

现在在四环算例上把两个方向逐步做一遍。原始图是环$1-2-3-4-1$。

\paragraph{由消元次序到袋树}
按次序$(1,2,3,4)$执行算法~\ref{alg:graph-dp-order-to-decomposition}，表~\ref{tab:graph-dp-four-cycle-construction}列出每一轮的状态。消去$1$时补入填充边$\{2,4\}$；消去$2$时当前邻居为$\{3,4\}$，已经相邻；之后不再补边。

\begin{table}[!htb]
  \centering
  \caption{四环上按次序$(1,2,3,4)$构建袋树。“当前邻居”从上一轮更新后的图读取；父袋按“后续邻居中次序最早者”确定，最后一列是与父袋的接口，它恰好等于当前邻居集合。}
  \label{tab:graph-dp-four-cycle-construction}
  \small
  \begin{tabular}{@{}ccccccc@{}}
    \toprule
    轮 & 消去 & 当前邻居$S_i$ & 填充边 & 消元簇$B_i$ & 父袋 & 接口\\
    \midrule
    $1$ & $1$ & $\{2,4\}$ & $\{2,4\}$ & $\{1,2,4\}$ & $B_2$ & $\{2,4\}$\\
    $2$ & $2$ & $\{3,4\}$ & 无 & $\{2,3,4\}$ & $B_3$ & $\{3,4\}$\\
    $3$ & $3$ & $\{4\}$ & 无 & $\{3,4\}$ & $B_4$ & $\{4\}$\\
    $4$ & $4$ & $\varnothing$ & 无 & $\{4\}$ & 根 & —\\
    \bottomrule
  \end{tabular}
\end{table}

得到的袋树是一条路径$\{1,2,4\}-\{2,3,4\}-\{3,4\}-\{4\}$，宽度为$2$。后两个袋包含于$\{2,3,4\}$，把它们并入后只剩$\{1,2,4\}-\{2,3,4\}$，接口为$\{2,4\}$，这就是第~\ref{sec:graph-dp-four-cycle}小节计算时使用的两个子问题。可以逐条核对定义：四条环边$\{1,2\},\{1,4\}$在第一个袋，$\{2,3\},\{3,4\}$在第二个袋；含$2$或$4$的袋相邻，含$1$或$3$的袋只有一个。

\paragraph{由袋树到消元次序}
反过来，从分解$\{1,2,4\}-\{2,3,4\}$出发执行算法~\ref{alg:graph-dp-decomposition-to-order}。以$\{2,3,4\}$为根时，叶袋$\{1,2,4\}$中不在父袋里的变量是$1$，先消去$1$，再在根袋中按$2,3,4$消去，读回次序$(1,2,3,4)$。

以$\{1,2,4\}$为根时，叶袋变成$\{2,3,4\}$，不在父袋里的变量是$3$。消去$3$时它在环上的当前邻居是$\{2,4\}$，补入$\{2,4\}$，这条填充边落在叶袋里；随后在根袋中按$1,2,4$消去，得到次序$(3,1,2,4)$，各步剩余邻居数为$2,2,1,0$。两个次序补入的是同一条填充边，诱导宽度都是$2$。

\paragraph{两种表示的对照}
两种表示中，最多同时处理的变量都是三个：消元簇$\{1,2,4\}$、$\{2,3,4\}$就是两个袋。由定理~\ref{thm:graph-dp-treewidth-induced-width}和第~\ref{sec:graph-dp-clique-containment}小节的下界，四环的树宽为$2$，所以四环上任何消元次序的诱导宽度都至少是$2$，这里的两个次序已经最优。

另一种同样宽度的分解是$\{1,2,3\}-\{1,3,4\}$，接口为$\{1,3\}$，它对应先消去$2$或$4$、补入$\{1,3\}$的次序。第~\ref{sec:graph-dp-query-structure}小节的保留集$R=\{1,3\}$要求的正是这种结构：在补边$\{1,3\}$后的图上只消去$V\setminus R$，按次序$(2,4)$执行算法~\ref{alg:graph-dp-order-to-decomposition}，消去$2$时$S_2=\{1,3\}\subseteq R$，消去$4$时$S_4=\{1,3\}\subseteq R$，两个消元簇$\{1,2,3\}$与$\{1,3,4\}$都连到根袋$B_R=\{1,3\}$。第~\ref{sec:graph-dp-retained-check}小节在这个结构上计算条件最小成本表。
